\chapter{Statistical physics on graphs}
\label{ch:introduction}
\index{branching process}

Statistical physics is the natural framework for problems in and on networks
\citep{martin2001,dorogovtsev2008}. It was built for systems with a great many
interacting parts, where the question is never what any one part does but what
the whole does on average --- and that is exactly the question a network poses.
A degree distribution is an ensemble; a giant component is an order parameter; a
percolation threshold is a phase transition. The correspondence is an identity
rather than an analogy, and it means that a century of machinery built for
magnets applies, with almost no translation, to questions
about connectivity, spreading and constraint satisfaction.

In this chapter we assemble that machinery: what a generating function does to a
degree distribution, what the cavity method is, and why replica symmetry
sometimes fails. None of it is new, and a reader who has it already should skip
to Chapter~\ref{ch:chygraphs} --- except for the last section, which is where
the book's own argument begins.

The argument runs as follows. Every method here works by pretending the network
is a tree. Real networks are not trees, and the specific
way they fail to be trees is that the neighbours of a node are neighbours of
one another. That is the
problem Chapter~\ref{ch:chygraphs} is designed to solve.

\section{Graphs and networks}
\label{sec:networks}
\index{degree distribution}
\index{excess degree|textbf}
\index{clustering coefficient|textbf}
\index{configuration model|textbf}
\index{generating function|textbf}

A \emph{network}, or graph, is a set of things together with a record of which
pairs of them are joined. The things are drawn as dots and called \emph{nodes};
a joining is drawn as a line and called a \emph{link}. That is the entire
vocabulary needed to start.

The \emph{degree} of a node is its number of links, written $k$. The
\emph{degree distribution} $p(k)$ is the fraction of nodes with exactly $k$
links --- equivalently, the probability that a node picked at random has degree
$k$. Its first moment $\ave{k}$ is twice the number of links divided by the
number of nodes, and its second moment $\ave{k^{2}}$ will turn out to matter
just as much.

Two more quantities, both used throughout.

\paragraph{The excess degree.} Pick a link at random and follow it. It ends at
a node, and the question that matters for everything downstream is not how many
links that node has but how many \emph{others} --- how many ways there are to
keep going. Because a node with $k$ links is $k$ times more likely to be found
at the end of a randomly chosen link than a node with one, the degree reached is
not drawn from $p(k)$ but from the size-biased distribution
$k\,p(k)/\ave{k}$, and the number of onward links is one less than that. Its
mean is the \emph{excess degree}
\begin{equation}
  \ave{\bar k}=\frac{\ave{k(k-1)}}{\ave{k}}
              =\frac{\ave{k^{2}}}{\ave{k}}-1.
  \label{eq:excessdeg}
\end{equation}
This is the single most important number in the first half of the book. It is
the average number of ways to continue after arriving somewhere, and almost
every threshold in what follows is the statement that this number, discounted by
something, equals one.

That it involves $\ave{k^{2}}$ rather than $\ave{k}$ is why heavy-tailed
networks behave so differently from ordinary ones. A network in which
$\ave{k^{2}}$ diverges has $\ave{\bar k}=\infty$, and thresholds that are
positive for a Poisson network go to zero for it.

\paragraph{The clustering coefficient.} Of all the pairs of a node's neighbours,
what fraction are themselves joined? If node $i$ has $k_{i}$ neighbours and
$e_{i}$ of the $\binom{k_{i}}{2}$ possible links among them exist, then
\begin{equation}
  C_{i}=\frac{2e_{i}}{k_{i}(k_{i}-1)},
  \label{eq:clustering}
\end{equation}
and $C$ is the average of $C_{i}$ over nodes of degree two or more. A high $C$
means the network is full of triangles. Table~\ref{tab:real} in
Chapter~\ref{ch:data} measures it for ten real networks; the values run from
$0.006$ to $0.65$. Section~\ref{sec:treebreaks} is about this quantity, and so,
in the end, is the book.

\paragraph{The configuration model.} Throughout, ``a random network with degree
distribution $p(k)$'' means the following construction. Give node $i$ a number
$k_{i}$ of stubs drawn from $p(k)$, then pair the stubs at random. What comes
out has the prescribed degree distribution and, in the limit of many nodes, no
other structure whatever: no clustering, no correlation between the degrees of
adjacent nodes, and --- the property that makes it tractable --- no short loops.
It is the null model against which everything else in network science is
measured, and it is also, not coincidentally, the object on which all the
methods of this chapter are exact.

\begin{figure}[t]
\centering
\begin{adjustbox}{max width=\linewidth}
\begin{tikzpicture}
\begin{scope}
  \node[tb,anchor=west] at (-0.3,2.75) {(a) degree};
  \node[hub] (H) at (1.25,1.35) {};
  \node[nd] (h1) at (0.15,2.05) {}; \node[nd] (h2) at (0.05,1.05) {};
  \node[nd] (h3) at (0.85,0.35) {}; \node[nd] (h4) at (1.85,2.15) {};
  \node[nd] (h5) at (2.35,1.15) {};
  \draw[lnk] (H)--(h1) (H)--(h2) (H)--(h3) (H)--(h4) (H)--(h5);
  \node[lb,anchor=west] at (1.55,1.35) {$k=5$};
  \node[lb,anchor=west,align=left] at (-0.3,-0.25)
    {how many links\\ a node has};
\end{scope}
\begin{scope}[xshift=4.35cm]
  \node[tb,anchor=west] at (-0.3,2.75) {(b) excess degree};
  \node[nd] (A) at (0.15,0.75) {};
  \node[ndf] (B) at (1.35,1.45) {};
  \node[nd] (b1) at (2.45,2.15) {}; \node[nd] (b2) at (2.55,1.25) {};
  \node[nd] (b3) at (2.15,0.45) {};
  \draw[lnk,line width=1.1pt] (A)--(B);
  \draw[lnk] (B)--(b1) (B)--(b2) (B)--(b3);
  \node[lb,anchor=south,rotate=30] at (0.72,0.98) {arrive};
  \node[lb,anchor=west] at (2.15,1.45) {$\bar k=3$};
  \node[lb,anchor=west,align=left] at (-0.3,-0.25)
    {how many ways there are\\ to keep going};
\end{scope}
\begin{scope}[xshift=9.4cm]
  \node[tb,anchor=west] at (-0.3,2.75) {(c) clustering};
  \node[nd] (C) at (1.15,1.35) {};
  \node[nd] (c1) at (0.15,2.05) {}; \node[nd] (c2) at (2.15,2.05) {};
  \node[nd] (c3) at (0.35,0.45) {}; \node[nd] (c4) at (2.05,0.45) {};
  \draw[lnk] (C)--(c1) (C)--(c2) (C)--(c3) (C)--(c4);
  \draw[lnk,line width=1.1pt] (c1)--(c2) (c3)--(c4);
  \node[lb,anchor=west,align=left] at (-0.3,-0.25)
    {of six neighbour pairs,\\ two are joined: $C_i=1/3$};
\end{scope}
\end{tikzpicture}
\end{adjustbox}
\caption{\textbf{Three quantities.} (a) Degree, the number of links at a node.
(b) Excess degree: arriving along a link, the number of ways on is one less than
the degree of the node arrived at, and that degree is size-biased because
well-connected nodes are found at the ends of more links. (c) Clustering, the
fraction of a node's neighbour pairs that are themselves joined --- the
quantity every method in this chapter assumes is zero, and no real network has
zero.}
\label{fig:basics}
\end{figure}

Figure~\ref{fig:basics} fixes those three quantities, and the third of them
is the one this chapter is about.

\section{Percolation}
\label{sec:percolation-intro}
\index{percolation!on a graph}
\index{Molloy--Reed criterion|textbf}
\index{giant component}

Take a network and dilute it. Keep each link with probability $p$ and delete it
otherwise. For small $p$ what remains is a scatter of small fragments.
For large $p$ a single component contains a finite fraction of all the nodes.
Somewhere between the two the fragments merge into that giant component, and
that abrupt change is the \emph{percolation transition}.

The question is where, and the answer follows from a counting argument that the
whole book generalises.

Start at a node and follow a surviving link. It ends somewhere, and from there
the exploration continues along $\bar k$ further links, of which a fraction $p$
survive. So one surviving branch begets, on average, $p\ave{\bar k}$ further
surviving branches. If that number is less than one the exploration dies out and
components stay finite; if it is greater than one it grows without bound and a
giant component exists. The threshold is
\begin{equation}
  p_{c}\ave{\bar k}=1,
  \qquad\text{that is,}\qquad
  p_{c}=\frac{\ave{k}}{\ave{k^{2}}-\ave{k}},
  \label{eq:molloyreed}
\end{equation}
which is the Molloy--Reed criterion \citep{molloy1995}. The generating-function
apparatus behind it is due to \citet{newman2001random}. For a Poisson network,
$\ave{\bar k}=\ave{k}$ and the giant component appears at mean degree one. For a
scale-free network with $\ave{k^{2}}$ divergent, $p_{c}=0$: such networks cannot
be broken by random deletion at all.

Three features of this argument carry through to
Chapter~\ref{ch:percolation} unchanged.

First, it is a \emph{branching process}. The exploration is treated as a family
tree in which each individual has an independent number of children, and the
threshold is the point where the mean number of children reaches one. This is
exact when the network has no short loops and wrong when it does, because on a
loop the exploration comes back to somewhere it has already been and a branching
process has no way to know that.

Second, the only property of the network that entered was $\ave{\bar k}$, a
ratio of the first two moments. Everything else about $p(k)$ is irrelevant to
where the transition is. It is not irrelevant to how big the giant component is
once past it, which is the subject of Chapter~\ref{ch:giant}.

Third, deleting links at random and deleting nodes at random are the same
calculation with different bookkeeping.

\begin{calculation}{the giant component, in full}
The branching argument gives the threshold. Getting the size of the giant
component needs generating functions, which are the workhorse of
Chapter~\ref{ch:percolation} and are introduced here once.

Encode the degree distribution as a power series,
\[
  G_{0}(x)=\sum_{k}p(k)\,x^{k},
\]
so that $G_{0}(1)=1$ and $G_{0}'(1)=\ave{k}$. The corresponding function for the
excess degree --- the number of onward links from a node arrived at along a link
--- is
\[
  G_{1}(x)=\frac{G_{0}'(x)}{G_{0}'(1)},
  \qquad G_{1}'(1)=\ave{\bar k}.
\]
Let $u$ be the probability that a link, followed, leads to a finite component.
The node at the far end must have all of \emph{its} onward links lead to finite
components, and each does so either because it was deleted (probability $1-p$)
or because it was kept and led to a finite component (probability $pu$). Summing
over the number of onward links,
\[
  u=G_{1}\bigl(1-p+pu\bigr).
\]
The fraction of nodes \emph{not} in the giant component is the same statement
made at a node rather than along a link, so the giant component fraction is
\[
  S=1-G_{0}\bigl(1-p+pu\bigr).
\]
The trivial solution $u=1$ always exists and gives $S=0$. It becomes unstable
when the derivative of the right-hand side at $u=1$ exceeds one, that is when
$p\,G_{1}'(1)>1$ --- Eq.~\eqref{eq:molloyreed} again. Threshold and order
parameter are the linearisation and the full solution of one equation, which is
the observation Chapter~\ref{ch:giant} is built on.
\end{calculation}

\begin{figure}[t]
\centering
\begin{adjustbox}{max width=0.95\linewidth}
\begin{tikzpicture}
% ---- branching picture
\node[ndf] (r) at (0,0) {};
\node[nd] (a1) at (1.5,1.1) {}; \node[nd] (a2) at (1.5,-1.1) {};
\node[nd] (b1) at (3.0,1.8) {}; \node[nd] (b2) at (3.0,0.7) {};
\node[nd] (b3) at (3.0,-0.5) {}; \node[nd] (b4) at (3.0,-1.7) {};
\node[nd] (c1) at (4.5,2.2) {}; \node[nd] (c2) at (4.5,1.4) {};
\node[nd] (c3) at (4.5,0.4) {}; \node[nd] (c4) at (4.5,-1.4) {};
\node[nd] (c5) at (4.5,-2.1) {};
\draw[lnk] (r)--(a1) (r)--(a2);
\draw[lnk] (a1)--(b1) (a1)--(b2) (a2)--(b3) (a2)--(b4);
\draw[lnk] (b1)--(c1) (b1)--(c2) (b2)--(c3) (b4)--(c4) (b4)--(c5);
\node[lb,anchor=south] at (0.75,0.60) {$1$};
\node[lb,anchor=west] at (5.0,1.7)
  {$\times\,p\ave{\bar k}$ per generation};
\draw[msg] (5.0,0.4) -- (6.1,0.4);
\node[lb,anchor=west,align=left] at (6.25,0.4)
  {$p\ave{\bar k}<1$: dies out\\
   $p\ave{\bar k}>1$: giant\\ \phantom{$p\ave{\bar k}>1$: }component};
\node[lb,anchor=north,align=center] at (2.2,-2.65)
  {exact on a tree, because no branch ever meets another};
\end{tikzpicture}
\end{adjustbox}
\caption{\textbf{Percolation as a branching process.} Following a surviving link
leads to a node with $\bar k$ further links, a fraction $p$ of which survive, so
each branch begets $p\ave{\bar k}$ branches. The transition is where that number
passes one. Every threshold in this book is a version of this statement; what
changes from chapter to chapter is what a branch is and what it carries.}
\label{fig:branching}
\end{figure}

Figure~\ref{fig:branching} is that argument drawn: one branch begetting more
than one branch, and why it is exact on a tree.

\section{Models on networks}
\label{sec:models}

Percolation asks about the network itself. The rest of the book puts something
\emph{on} the network --- a variable at every node, and an energy that says which
combinations of variables are cheap. Three of the models that run through the book are introduced here, and all
three appear now on the same small graph, so the family resemblance is
visible. For the wider map --- spin models, $k$-core percolation, epidemic
thresholds and the rest, on equilibrium and growing networks alike ---
\citet{dorogovtsev2008} is the review to read alongside this chapter, which
takes only what the later chapters use.

\subsection{Ising model}
\label{sec:ising-intro}
\index{Ising model|textbf}

Put a variable $\sigma_{i}=\pm1$ on every node --- historically a magnetic
moment pointing up or down --- and let neighbouring nodes prefer to agree:
\begin{equation}
  \mathcal{H}=-J\sum_{\langle ij\rangle}\sigma_{i}\sigma_{j}
              -B\sum_{i}\sigma_{i},
  \label{eq:ising-intro}
\end{equation}
the first sum over links, the second over nodes, $J>0$ the strength of the
preference and $B$ an external field pushing every spin one way. Configurations
occur with the Boltzmann probability $\propto e^{-\beta\mathcal{H}}$, where
$\beta=1/T$ is the inverse temperature.

The behaviour is a competition. At high temperature the spins are essentially
independent and the average magnetisation $m=\ave{\sigma}$ is zero. At low
temperature agreement wins, a majority forms, and $m\ne0$ even with $B=0$. In
between is a critical temperature $T_{c}$.

On a treelike network the answer is a one-line calculation, and it is
\emph{the same line as percolation}. Write $t=\tanh\beta J$ for the strength
with which one spin transmits its bias to a neighbour. Then
\begin{equation}
  t\,\ave{\bar k}=1
  \label{eq:tc-intro}
\end{equation}
locates $T_{c}$. Compare Eq.~\eqref{eq:molloyreed}: the occupation probability
$p$ has been replaced by $\tanh\beta J$ and nothing else has changed. A
correlation propagating through the network and a cluster spreading through it
are the same branching process counted with a different weight.
Chapter~\ref{ch:potts} shows that the two are the same model at two values of
one parameter.

Equation~\eqref{eq:tc-intro} is exact on a random graph with any degree
distribution, and the same replica calculation returns the critical exponents
from the moments of $p(k)$ \citep{leone2002}. The parallel with percolation
survives into the pathologies: where $\ave{k^{2}}$ diverges $\ave{\bar k}$ is
infinite, Eq.~\eqref{eq:tc-intro} has no finite root and the system is ordered
at every temperature, which is the counterpart of $p_{c}=0$ above; where
$\ave{k^{2}}$ is finite but $\ave{k^{4}}$ is not, $T_{c}$ is finite and the
exponents are not the mean-field ones.

\subsection{Vertex covering}
\label{sec:cover-intro}
\index{vertex cover|textbf}
\index{NP-hardness}

The second model has no temperature in it at all. A \emph{vertex cover} of a graph
is a set of nodes such that every link has at least one endpoint in the set ---
guards posted at corridor junctions so that every corridor is watched. Covering
everything is easy; covering everything with as few nodes as possible is the
\emph{minimum vertex cover} problem, and it is NP-hard \citep{karp1972}: no
algorithm is known that solves it in time polynomial in the size of the graph,
and it is widely believed none exists.

Statistical physics gets a grip on it by making it a model. Put
$x_{i}\in\{0,1\}$ on every node, $1$ meaning covered, and write
\begin{equation}
  \mathcal{H}=\sum_{i}x_{i},
  \qquad\text{subject to}\quad
  x_{i}+x_{j}\ge1 \ \ \text{on every link}.
  \label{eq:vc-intro}
\end{equation}
The energy is just the number of guards, the constraint is hard --- infinitely
costly to violate --- and the minimum vertex cover is the ground state. Asking
for the ground state of a disordered system is asking an optimisation question,
and the whole apparatus of statistical mechanics becomes available for it. This
was the move that opened the subject \citep{weigt2000}.

One fact about vertex cover will matter repeatedly. Consider a node of degree
one --- a \emph{leaf}. Its single link must be covered, and covering it at the
leaf can never be better than covering it at the neighbour, since the neighbour
may cover other links too. So there is always a minimum cover that takes the
neighbour and not the leaf, and one can simply remove both, together with every
link the neighbour touched, and repeat. This is \emph{leaf removal}
\citep{bauer2001}. When it consumes the whole graph it has found a minimum cover
and \emph{proved} it minimum. When it stalls, what is left is called the
\emph{core}, and it is the hard part.

For a Poisson random graph, leaf removal succeeds when the mean degree is below
$e=2.718\ldots$ and stalls above it \citep{bauer2001,weigt2000}. That number
will reappear, from a completely different direction, in
Chapter~\ref{ch:hittingset}.

\subsection{Hitting set}
\label{sec:hs-intro}
\index{hitting set problem|textbf}

Vertex cover generalises in the obvious direction. Instead of links --- sets of
two nodes --- take a family of sets of any size, and ask for the smallest
collection of elements meeting every one of them. That is the \emph{hitting set}
problem \citep{mezard2007}, and vertex cover is the special case in which every
set has two members.

The Hamiltonian is Eq.~\eqref{eq:vc-intro} with the constraint read over
hyperedges rather than links, $\sum_{i\in a}x_{i}\ge1$ for every set $a$.
Nothing about the formulation changes; what changes is that the constraint now
couples three, four or ten variables at once, and Chapter~\ref{ch:hittingset}
shows that a limit everybody takes for granted on graphs stops being correct the
moment it does.

This pairing --- vertex cover as the two-element case of hitting set --- is the
reason Chapters~\ref{ch:hittingset} and~\ref{ch:cover} are adjacent. They are one
problem.

\begin{figure}[t]
\centering
\begin{adjustbox}{max width=\linewidth}
\begin{tikzpicture}
\def\gA{\node[nd] (v1) at (0.15,1.55) {}; \node[nd] (v2) at (1.25,2.05) {};
        \node[nd] (v3) at (1.15,0.85) {}; \node[nd] (v4) at (2.35,1.55) {};
        \node[nd] (v5) at (2.25,0.35) {};}
% ---- Ising
\begin{scope}
  \node[tb,anchor=west] at (-0.3,2.85) {(a) Ising};
  \node[nd] (a1) at (0.15,1.55) {\tiny$+$}; \node[nd] (a2) at (1.25,2.05) {\tiny$+$};
  \node[nd] (a3) at (1.15,0.85) {\tiny$-$}; \node[nd] (a4) at (2.35,1.55) {\tiny$+$};
  \node[nd] (a5) at (2.25,0.35) {\tiny$-$};
  \draw[lnk] (a1)--(a2) (a1)--(a3) (a2)--(a3) (a2)--(a4) (a3)--(a5) (a4)--(a5);
  \node[lb,anchor=west,align=left] at (-0.3,-0.30)
    {a sign at every node;\\ neighbours prefer to agree};
\end{scope}
% ---- vertex cover
\begin{scope}[xshift=4.3cm]
  \node[tb,anchor=west] at (-0.3,2.85) {(b) vertex cover};
  \node[nd]  (b1) at (0.15,1.55) {}; \node[ndf] (b2) at (1.25,2.05) {};
  \node[ndf] (b3) at (1.15,0.85) {}; \node[nd]  (b4) at (2.35,1.55) {};
  \node[ndf] (b5) at (2.25,0.35) {};
  \draw[lnk] (b1)--(b2) (b1)--(b3) (b2)--(b3) (b2)--(b4) (b3)--(b5) (b4)--(b5);
  \node[lb,anchor=west,align=left] at (-0.3,-0.30)
    {a set of nodes touching\\ every link, as small as possible};
\end{scope}
% ---- hitting set
\begin{scope}[xshift=8.9cm]
  \node[tb,anchor=west] at (-0.3,2.85) {(c) hitting set};
  \node[nd]  (c1) at (0.15,2.15) {}; \node[nd]  (c2) at (0.15,1.05) {};
  \node[ndf] (c3) at (1.30,1.60) {}; \node[ndf] (c4) at (2.40,1.60) {};
  \node[nd]  (c5) at (3.35,0.85) {};
  \begin{scope}[on background layer]
    \node[cx,fill=black!26,fill opacity=0.7,fit=(c1)(c2)(c3)]{};
    \node[cx,fill=black!26,fill opacity=0.7,inner sep=3.8pt,fit=(c3)(c4)]{};
    \node[cx,fill=black!26,fill opacity=0.7,fit=(c4)(c5)]{};
  \end{scope}
  \node[lb,anchor=west,align=left] at (-0.3,-0.30)
    {a set of nodes touching\\ every \emph{group}, as small as possible};
\end{scope}
\end{tikzpicture}
\end{adjustbox}
\caption{\textbf{Three models, one graph.} (a) The Ising model decorates the
nodes with signs and pays for disagreement across links. (b) Vertex cover picks
a set of nodes (filled) meeting every link. (c) Hitting set is the same question
with links replaced by groups of any size --- and the groups, drawn here as
boxes, are already the complexes of Chapter~\ref{ch:chygraphs}. What the three
share is a variable per node, an energy, and a question about the resulting
Boltzmann measure; what they differ in is what the message of
Sec.~\ref{sec:cavity} has to carry.}
\label{fig:threemodels}
\end{figure}

Figure~\ref{fig:threemodels} puts the three side by side on one graph: they
differ in what decorates the nodes and what is paid for, not in the object
underneath.

\subsection{What these have in common}

All three are a variable at every node, an energy summing local terms, and a
Boltzmann measure $\propto e^{-\beta\mathcal{H}}$. Optimisation is the
zero-temperature limit of that measure: as $\beta\to\infty$ all the weight
concentrates on the lowest-energy configurations, so ``find the minimum vertex
cover'' and ``take $T\to0$ in this model'' are the same instruction. The
apparatus does not distinguish them, which is why one method can serve a
magnet and an NP-hard optimisation problem.

One caveat about that limit runs through the book. Taking
$T\to0$ inside an approximation is not the same as taking it in the exact
answer, and Chapter~\ref{ch:hittingset} exhibits a case, simple enough to check
by hand, where the standard zero-temperature treatment gives $1/2$ and the truth
is $1/3$.

\section{Mean field}
\label{sec:meanfield}
\index{mean-field theory|textbf}

The oldest way to solve a model like Eq.~\eqref{eq:ising-intro} is to assume
each spin feels not its neighbours but their average. Replace
$\sigma_{j}\to m$ for every neighbour, and a node of degree $k$ then sits in an
effective field $Jkm$, giving the self-consistency condition
\begin{equation}
  m=\tanh\bigl(\beta J\ave{k}m\bigr).
  \label{eq:mf}
\end{equation}
For small $m$ the right-hand side is $\beta J\ave{k}m$, so a non-zero solution
appears when $\beta J\ave{k}=1$.

This is wrong. Compare it with Eq.~\eqref{eq:tc-intro}: mean field gives $\beta J\ave{k}=1$ where the correct
treelike answer is $\tanh(\beta J)\ave{\bar k}=1$. Two errors, in opposite
directions. Mean field uses $\ave{k}$ where it should use the excess degree
$\ave{\bar k}$ --- it forgets that the walk arrived along a link and cannot
leave by it --- and it uses $\beta J$ where it should use $\tanh\beta J$,
because it treats the neighbour's spin as a fixed number rather than a
fluctuating variable that only partially transmits.

Mean field becomes exact when every node has many neighbours, since then the sum
of many weakly correlated influences really is close to its average. In infinite
dimension, or on a fully connected graph, it is right. On a sparse network,
where a node has three or five neighbours and each one matters individually, it
is not, and the fix is the subject of the next section.

What survives from mean field is the \emph{style} of answer: a self-consistency
condition for one number, solved by looking for where a non-trivial solution
appears. Everything in this book has that shape. What changes is that the one
number becomes a message, and then a distribution of messages.

\section{The cavity method}
\label{sec:cavity}
\index{cavity method|textbf}
\index{belief propagation|textbf}
\index{Bethe--Peierls approximation|textbf}
\index{message passing}

Our core idea is geometric before it is a calculation.

Take a node $i$ and remove it, leaving a hole --- a \emph{cavity}. On a tree,
removing $i$ disconnects its neighbours from one another completely: the
branches that hung off $i$ are now separate networks with no path between them.
Whatever is going on in one branch is therefore independent of whatever is going
on in another. That independence is the trick, and it converts a problem
in which every variable is coupled to every other into a problem about one
branch at a time.

Concretely, let each neighbour $j$ of $i$ send $i$ a summary of what its branch
would do if $i$ were absent. Call that summary a \emph{message}. For the Ising
model the message is a field $h_{j\to i}$: the effective bias that branch $j$
exerts on node $i$. Because the branches are independent, the total bias on $i$
is just the sum of the messages arriving at it, and the message $i$ sends onward
to some other neighbour $l$ is built from the messages $i$ receives from
everyone except $l$:
\begin{equation}
  h_{i\to l}=B+\sum_{j\in\partial i\setminus l}
             u\bigl(h_{j\to i}\bigr),
  \label{eq:cavity-intro}
\end{equation}
where $\partial i$ is the set of neighbours of $i$ and $u(\cdot)$ is the
function that says how much of a bias survives one link. This is the
\emph{cavity recursion}. Under other names it is the Bethe--Peierls
approximation \citep{bethe1935} and \emph{belief propagation}
\citep{yedidia2005}; they are the same equations arrived at by physicists,
statisticians and computer scientists independently; \citet{mezard2009} set the
three vocabularies side by side.

\begin{figure}[t]
\centering
\begin{adjustbox}{max width=0.95\linewidth}
\begin{tikzpicture}
% node i, removed: drawn as a dashed outline so the hole it leaves is visible
\node[circle,draw=black!55,dashed,fill=white,minimum size=3.6mm,inner sep=0pt,
      line width=.5pt] (i) at (0,0) {};
\node[lb,anchor=north] at (0,-0.30) {$i$, removed};
% its three neighbours, and the branch hanging off each
\node[nd] (j1) at (-1.75,1.05) {}; \node[nd] (j2) at (1.75,1.05) {};
\node[nd] (j3) at (0,-1.55) {};
\node[lb,anchor=south east] at (-1.85,1.15) {$j$};
\draw[lnk,dashed] (i)--(j1) (i)--(j2) (i)--(j3);
\node[nd] (p1) at (-3.05,1.95) {}; \node[nd] (p2) at (-1.60,2.30) {};
\node[nd] (p3) at (-3.45,0.85) {};
\draw[lnk] (j1)--(p1) (j1)--(p2) (p1)--(p3);
\node[nd] (q1) at (3.05,1.95) {}; \node[nd] (q2) at (1.60,2.30) {};
\node[nd] (q3) at (3.45,0.85) {};
\draw[lnk] (j2)--(q1) (j2)--(q2) (q1)--(q3);
\node[nd] (r1) at (-1.15,-2.35) {}; \node[nd] (r2) at (1.15,-2.35) {};
\draw[lnk] (j3)--(r1) (j3)--(r2);
% the message, with the label set clear of the arrow
\draw[msg] (-1.40,0.84) -- (-0.34,0.20);
\node[lb,anchor=east] at (-0.92,0.26) {$h_{j\to i}$};
% the three branches, named
\draw[decorate,decoration={brace,amplitude=3pt},black!45]
  (-3.75,2.55) -- (-1.30,2.55);
\node[lb,anchor=south] at (-2.52,2.68) {branch};
\draw[decorate,decoration={brace,amplitude=3pt},black!45]
  (1.30,2.55) -- (3.75,2.55);
\node[lb,anchor=south] at (2.52,2.68) {branch};
\draw[decorate,decoration={brace,amplitude=3pt,mirror},black!45]
  (-1.45,-2.72) -- (1.45,-2.72);
\node[lb,anchor=north] at (0,-2.88) {branch};
\end{tikzpicture}
\end{adjustbox}
\caption{\textbf{The cavity.} Removing node $i$ from a tree separates its
branches completely. Each branch then summarises itself in a single message
$h_{j\to i}$, and because the branches are independent those messages simply
add. Equation~\eqref{eq:cavity-intro} is that statement written down. Everything
in Chapters~\ref{ch:statmech} to~\ref{ch:satisfiability} is this picture with a
complex in place of a node, and the only question that ever changes is what the
message carries.}
\label{fig:cavity}
\end{figure}

Figure~\ref{fig:cavity} is the construction, and the whole of the method is in
it: remove a node, and what is left summarises itself.

Two ways of using Eq.~\eqref{eq:cavity-intro} run side by side throughout the
book and should not be confused.

\emph{On one instance}, one has an actual graph, initialises the messages
arbitrarily, and iterates until they stop changing. That is belief propagation,
it is an algorithm, and it returns an answer for that particular graph.

\emph{On an ensemble}, one gives up on individual messages and asks for their
\emph{distribution} over a random graph. Equation~\eqref{eq:cavity-intro} then
becomes an equation for a probability distribution $P(h)$, self-consistent in
the sense that a message built out of $\bar k$ messages drawn from $P$ is itself
drawn from $P$. This gives the typical behaviour of the whole class of graphs at
once, and it is what all the thresholds in this book are.

\begin{calculation}{the cavity recursion for the Ising model, and its threshold}
For Eq.~\eqref{eq:ising-intro} the message function is
\[
  u(h)=\frac{1}{\beta}\,\mathrm{atanh}
       \bigl[\tanh(\beta J)\tanh(\beta h)\bigr],
\]
which is derived by summing over the single spin $\sigma_{j}$ at the far end of
the link: a spin biased by $h$ transmits a bias reduced by the factor
$t=\tanh\beta J$, and the composition of two hyperbolic tangents is what
enforces that the transmitted bias can never exceed the link's own strength.

Set $B=0$ and ask when the paramagnetic solution $h\equiv0$ becomes unstable.
For small $h$,
\[
  u(h)\simeq t\,h,
\]
so Eq.~\eqref{eq:cavity-intro} linearises to
$h_{i\to l}=t\sum_{j\in\partial i\setminus l}h_{j\to i}$: a small bias is
multiplied by $t$ at every step and fans out over $\bar k$ branches. Averaging
over the ensemble, a perturbation grows when
\[
  t\,\ave{\bar k}>1,
\]
which is Eq.~\eqref{eq:tc-intro}. For a $c$-regular graph $\ave{\bar k}=c-1$ and
this is the classical Bethe-lattice result $\tanh\beta J_{c}=1/(c-1)$.

Three features of this calculation are what Chapter~\ref{ch:statmech}
generalises. The message is a \emph{number} here
only because the interaction is on a single link; when the interaction lives
inside a group, the message becomes a distribution. The factor $t$ came from
summing over the interior of one link \emph{exactly}; when the interior is
larger, that sum is larger but still exact. And the counting $\ave{\bar k}$ is
independent of both --- it is a property of the network, not of the model.
Separating those three is what lets one calculation serve every model in this
book.
\end{calculation}

\section{Replica symmetry and symmetry breaking}
\label{sec:replicas}
\index{replica symmetry breaking|textbf}
\index{de Almeida--Thouless line}

The cavity method as just described contains an assumption that is easy to miss.
In writing a single distribution $P(h)$ and asking it to be self-consistent, one
has assumed that the system has essentially \emph{one} state --- that two runs
of the model would look statistically alike. This is
\emph{replica symmetry}, so called because in the older replica formalism
\citep{mezard1987} it appears as a symmetry among $n$ copies of the system.

For the Ising model at high temperature it is true. For an Ising model with
random signs of $J$, a spin glass, it is false below a certain temperature: the
energy landscape breaks into exponentially many valleys separated by barriers,
each valley a perfectly good ``state'', and no single $P(h)$ describes them all.
The same happens in hard optimisation problems --- for vertex cover on a Poisson
graph it happens at mean degree $e$, the point where leaf removal stalls
\citep{weigt2000,bauer2001}; Chapter~\ref{ch:cover} shows that the two are one
point.

The remedy is \emph{one-step replica symmetry breaking}: a distribution over
valleys, and within each valley a distribution over messages --- a distribution
of distributions \citep{mezard2001}. This book does that calculation twice, in
Secs.~\ref{sec:colouring-sp} and~\ref{sec:sat-sp}, and only where the messages
are warnings and the distribution is over a finite alphabet; everywhere else it
does not. So it needs to know when the calculation it \emph{does} do has
stopped being right, and there are three reliable signals.

\emph{The stability of the symmetric solution.} Perturb the messages and ask
whether the perturbation grows. For the Ising model the relevant condition turns
out to be $t^{2}\ave{\bar k}=1$ rather than $t\ave{\bar k}=1$ --- the
\emph{de Almeida--Thouless line} \citep{dealmeida1978} --- because a perturbation
of \emph{random sign} is transmitted with $t$ and its effect measured by
squaring. That squaring is the difference between the ferromagnetic
transition and the spin-glass one, and Chapter~\ref{ch:ising} obtains both from
the same expression by squaring one factor.

\emph{Negative entropy.} A replica-symmetric calculation of a discrete problem
returns the logarithm of a number of configurations. Where that goes negative
the calculation is counting fewer than one configuration, which is impossible,
and the answer it gives is an underestimate.

\emph{The iteration stops converging.} Run the equations and they cycle or
wander. This is the practical signal and it is the one that usually arrives
first.

None of these gives the right answer. All of them show that the answer in hand
is wrong. This book flags every place they fire.

\section{The tree assumption, and why real networks break it}
\label{sec:treebreaks}
\index{tree assumption|textbf}
\index{treelike}
\index{clustering!and the tree assumption}

In this section we collect the tree assumption made by every method above and
measure how badly real networks violate it, since that violation is what the
rest of the book is written to repair.

The percolation threshold came from a
branching process, which is exact when no branch ever meets another. The cavity
recursion came from removing a node and finding its branches independent, which
is true when the branches are not joined behind its back. Both statements are
the same statement: \textbf{the network must be locally treelike}. Follow links
away from a node must not lead back to it.

Real networks come back. The clustering coefficient of
Eq.~\eqref{eq:clustering} measures exactly how often, and it is not small: in
the ten networks of Table~\ref{tab:real} it runs from $0.006$ to $0.65$. A
triangle is a loop of length three, the shortest possible, and every triangle is
a place where two branches that the theory believes to be independent are in
fact the same branch.

\begin{figure}[t]
\centering
\begin{adjustbox}{max width=0.95\linewidth}
\begin{tikzpicture}
\begin{scope}
  \node[tb,anchor=west] at (-0.3,2.35) {(a) what the theory assumes};
  \node[ndf] (i) at (1.1,1.55) {};
  \node[nd] (j) at (0.15,0.65) {}; \node[nd] (l) at (2.05,0.65) {};
  \node[nd] (j1) at (-0.55,-0.35) {}; \node[nd] (l1) at (2.75,-0.35) {};
  \draw[lnk] (i)--(j) (i)--(l) (j)--(j1) (l)--(l1);
  \node[lb,anchor=west,align=left] at (-0.3,-1.05)
    {two branches, never\\ meeting: independent};
\end{scope}
\begin{scope}[xshift=5.3cm]
  \node[tb,anchor=west] at (-0.3,2.35) {(b) what is there};
  \node[ndf] (I) at (1.1,1.55) {};
  \node[nd] (J) at (0.15,0.65) {}; \node[nd] (L) at (2.05,0.65) {};
  \node[nd] (J1) at (-0.55,-0.35) {}; \node[nd] (L1) at (2.75,-0.35) {};
  \draw[lnk] (I)--(J) (I)--(L) (J)--(J1) (L)--(L1);
  \draw[lnk,line width=1.2pt] (J)--(L);
  \node[lb,anchor=west,align=left] at (-0.3,-1.05)
    {one link closes a triangle,\\ and they are not};
\end{scope}
\end{tikzpicture}
\end{adjustbox}
\caption{\textbf{The assumption, and the world.} The cavity recursion adds the
messages arriving at a node because the branches sending them cannot influence
one another. One extra link --- the heavy one --- makes that false, and there is
no repair within the graph: the two messages are now correlated in a way a sum
cannot express. Chapter~\ref{ch:chygraphs} does not repair it either. It stops
treating the triangle as three links.}
\label{fig:treebreaks}
\end{figure}

Figure~\ref{fig:treebreaks} draws the gap between the two.

How much does this matter? Two answers, and they are of different kinds.

\paragraph{Quantitatively, it is a wrong number.} Take a network with two
triangles at every node and compare it with a $4$-regular network --- same
degree everywhere, same $\ave{\bar k}$, so Eq.~\eqref{eq:tc-intro} gives them the
same critical temperature. It is wrong by $13.9$ per cent, and
Chapter~\ref{ch:ising} computes the correct value in closed form and confirms it
by Monte Carlo at $14.2$ per cent. The sign is not the one most people guess:
clustering \emph{lowers} $T_{c}$.

\paragraph{Qualitatively, it is the wrong phase.} Hyperbolic random graphs, which
are clustered because of the geometry they are built in, retain an extensive
leaf-removal core at \emph{every} mean degree, with no transition at all. Their
degree-matched configuration model --- the same $p(k)$, the same $\ave{\bar k}$,
the same assortativity to within a few per cent --- has no core whatever until
the mean degree reaches somewhere between $4.5$ and $10$. The two graphs agree
on every quantity the theory of this chapter can see, and they are in different
phases. Whatever separates them lies outside $p(k)$ by construction, and it is
clustering.

The second example rules out the obvious response. One cannot rescue the situation by measuring the degree distribution
more carefully, or by correlating the degrees of adjacent nodes
\citep{vazquez2003vw}, or by any refinement of a description whose atoms are
nodes and links. The
information that distinguishes the two graphs is not in that description at all.

So the field has had, for twenty-five years, a method that works and a
systematic reason why it should not apply to the objects it is applied to.

\section{Generalised belief propagation}
\label{sec:gbp-intro}
\index{generalised belief propagation}
\index{Kikuchi approximation}
\index{cluster variation method}
\index{region graph}

There is an established answer to loops, and it is worth stating here because
this book uses it later and because it is the alternative to the route the rest
of the book takes.

Belief propagation is exact on a tree because a tree decomposes into nodes and
links with every variable, and every interaction, counted exactly once. When the
network is not a tree that bookkeeping fails. The remedy, due to Kikuchi and
given its message-passing form by \citet{yedidia2005}, is to keep the
bookkeeping and enlarge the units: work with \emph{regions} --- sets of
variables treated as single objects --- and close the family of regions under
intersection. Each region $R$ is assigned a counting number fixed by
M\"obius inversion,
\begin{equation}
  c_{R}=1-\sum_{R'\supsetneq R}c_{R'},
  \label{eq:mobius-intro}
\end{equation}
which is precisely the condition that every variable and every interaction is
counted once across the whole family. Messages then pass from each region to its
sub-regions rather than from node to node.

Three properties matter for what follows. The scheme \emph{contains} the
cavity method: when the regions are the interactions and the single variables,
Eq.~\eqref{eq:mobius-intro} returns $c_{v}=1-k$ and the update reduces to
Eq.~\eqref{eq:cavity-intro}. It is \emph{exact} whenever the closed family of
regions forms a junction tree, which is the loop-free condition one level up. And
it costs $2^{|R|}$ per region, so it buys accuracy with enumeration, and the
purchase stops being affordable when the regions grow.

What it does not give is an ensemble. The regions are an explicit list drawn
from one network, and the messages are functions on them, where everything else
in this chapter carries a distribution over a random graph.
Chapter~\ref{ch:gbp} measures what this method recovers, and what it costs, on
three classes of network.

The next chapter's response is not to patch the method. It is to change what the
method is applied to: if a triangle is a problem when treated as three links,
stop treating it as three links.

